\section{Method}
\label{sec:method}
AIMI uses in-hand motion to acquire visual evidence and turn it into a presentation goal. An agent decides how to inspect the object, while a recurrent motor policy and local controllers maintain continuous control. Training proceeds from a ViserDex-based teacher--student foundation adapted to Revo3, to goal-bank adaptation, and then local-controller learning (Fig.~\ref{fig:pipeline}). We describe the cube implementation below. Other objects use their own students and maintenance controllers; their interfaces are given in the supplement.

\section{Task and Evaluation}
\label{sec:task}
\paragraph{Semantic target.}
Q1 evaluates six objects in simulation: a marked cube, a wooden die, a duck, a hedgehog, a pig, and a milk carton. Q2 and Q3 use only the marked cube. Instructions request one of six surface letters, or one through six dots on the die. The cube letters are $\{T,E,R,L,S,H\}$. Each target has a fixed reference orientation $Q(g)$ that places the requested face upward. Scoring compares the full quaternion, including rotation about the face normal. This rule is shared by all methods.

At step $t$, the agent receives an RGB image $I_t$ and instruction $\ell$. The motor policy receives proprioception $x_t$ and measured pose $\hat p_t$. Cube pose is estimated from RGB; other objects use noisy measurements constructed from simulator pose. Controllers do not directly receive the unperturbed true pose $p_t$. Teacher actions and privileged targets are used for training, and true pose is used for independent scoring. During deployment, the controller cannot use a future agent reply. We denote the active goal vector by $g_t$, the recurrent state by $h_t$, and a pending request by $m_t=1$.

\paragraph{Timing.}
Only one agent request can be pending at a time. Cube episodes allow up to six model calls and 3,600 control steps at 30~Hz. Other objects use 20~Hz over the same 120 seconds of simulated time. For the cube, a call taking $L$ wall-clock seconds can affect control only after it completes and at least $\lceil30L\rceil$ steps have passed since the request. Physics continues while the agent reasons. The simulator runs slower than real time, so this models inference delay but does not demonstrate a 30~Hz hardware system.

\paragraph{Autonomous completion.}
All cube systems use the same stopping check at 3~Hz. The agent must have bound the requested symbol to a goal, with no request pending and no inspection underway. Measured orientation error must be at most $\epsilon=0.16$ radians. This check can stop the system without an explicit stop action from the agent. On other objects, REVEAL uses calibrated completion correction, while the baselines use a measured-pose check. Their stopping rules are detailed in the supplement, and all are scored with the same true-pose tolerance. At stop time $\tau$, success is
\begin{equation}
 S=\mathbf{1}\!\left[\hat b_\tau=b^*(\ell),\ d(q_\tau,Q(b^*))\leq\epsilon,\ \neg D_{0:\tau}\right],
\label{eq:success}
\end{equation}
where $\hat b_\tau$ is the declared target binding, $q_\tau$ is true orientation, $D_{0:\tau}$ indicates an earlier terminal drop, and $d(q,Q)=2\arccos(|\langle q,Q\rangle|)$ for normalized quaternions. The target must be correct at the stop time, not merely reached earlier. We do not require an additional hold after stopping. A wrong binding or an incorrect final pose counts as failure even if the object remains in the hand. A drop terminates the trial when the object is at least 0.24 metres from its nominal grasp position.

\paragraph{Paired trials and service failures.}
Unless stated otherwise, each condition uses 40 trials: 20 per seed under two seeds. Compared methods share the same initial physical states and scored targets. Control and feedback experiments also share instructions. In the mapping experiment, the requested letter changes with the layout so that the physical target stays the same. A service error or a model call taking more than 120 seconds excludes that trial from the service-eligible comparison. We remove the same trial from every compared method, retain all raw results, and do not replace excluded trials. Each reported rate uses the denominator for its own comparison. Service failures are reported separately from control failures.


\subsection{Learning goal-conditioned motor skills}
\paragraph{Teacher--student foundation.}
We build on the ViserDex teacher--student approach adapted to the Revo3 hand~\cite{bhardwaj2026viserdex}. The pretrained teacher uses privileged simulation information and geometric goals to provide action supervision. Starting from the corresponding student, we retain compatible encoders, recurrent and actor weights, and normalization statistics. The teacher remains frozen during the following adaptation.

\paragraph{Goal-bank adaptation.}
To connect discrete agent decisions to motor skills, we use six normalized vectors $B=\{g_1,\ldots,g_6\}$ from cached SigLIP2 text features~\cite{tschannen2025siglip2}. Each represents a color skill and its reference object orientation. A learned projection replaces the student's geometric goal inputs:
\begin{align}
 z_t&=\tanh(W_g g_t+b_g),\\
 (u_t,h_{t+1})&=\pi_\theta([x_t,z_t],\hat p_t,h_t).
\label{eq:student}
\end{align}
The adapted student receives neither the target quaternion nor the relative-goal quaternion. The goal bank is fixed; online agent outputs select its entries rather than supplying hidden model features.

We fit the projection and recurrent policy to frozen-teacher actions using a DAgger-style procedure~\cite{ross2011dagger}. Projection warmup uses action MSE. Recurrent adaptation adds a training-only decoder that reconstructs normalized privileged state with the goal quaternion removed. Its reconstruction MSE has weight 0.2. The teacher and decoder are removed from deployment. Network dimensions and fitting details are provided in the supplement.

\subsection{Agentic visual exploration}
\label{sec:exploration}
The agent receives the instruction, current image, executed-action history and measured-state summaries, and returns an action: inspect, reorient, verify, stop, or abstain. State summaries include estimated skill-relative errors and current binding/request status, but no true pose or correctness label. We validate replies before applying them. Inspection selects a skill to obtain another view; target binding selects the motor goal for presentation.


The agent chooses inspection and reorientation durations of 10--90 control steps. Later images reveal the effect of these actions and support the next decision. Action duration is distinct from model latency. Only one request can be pending, and the current skill remains active until a valid reply changes it. Invalid replies never select an oracle target. This loop lets the hand gather evidence before committing to a presentation goal, while control continues between decisions.

\subsection{Continuous control during and after reasoning}
\label{sec:wait}
The agent chooses goals at irregular intervals, but the hand needs commands at every control step. We therefore place local controllers around the adapted student. A sensor encoder and a goal encoder feed a fusion layer that predicts command corrections. Sensor features include proprioception, measured pose, the student's recurrent state and proposal, pose-fitting error, and binding and request status.

\paragraph{Waiting and execution commands.}
The waiting actor also receives the pose and command recorded when the request began, called the \emph{request anchor}, and elapsed waiting time. It predicts a speed logit $r_t$ and joint correction values $v_t$. With normalized current joints $\bar q_t$ and clipped student targets $\bar u_t$, the command is
\begin{align}
 \alpha_t&=\sigmoid(4r_t+\logit(0.15)),\\
 b_t&=\begin{cases}
 \bar q_t+\alpha_t(\bar u_t-\bar q_t),&m_t=1,\\
 \bar u_t,&m_t=0,
 \end{cases}\\
 a_t&=\clip\bigl(b_t+0.35\tanh(v_t),-1,1\bigr).
\label{eq:control}
\end{align}
When no request is pending, a separate execution actor supplies $v_t$ at full proposal speed. Zero waiting-actor outputs recover fixed slowing with coefficient 0.15. Joint correction can also change command direction, which scalar slowing alone cannot do. The underlying joint mapping and smoothing remain unchanged.


\paragraph{Teacher-supervised correction.}
The frozen teacher supplies holding actions for the request anchor during waiting and actions toward the active goal during execution. We train each actor against these labels with a penalty on its correction $\delta_\phi=0.35\tanh(v_t)$:
\begin{equation}
 \mathcal{L}_{\mathrm{control}}=\mathbb{E}_{(s,g,a^*)\sim\mathcal{D}}\left[\|a_\phi(s,g)-a^*\|_2^2+0.01\|\delta_\phi(s,g)\|_2^2\right].
\label{eq:control_loss}
\end{equation}
Here $s$ collects sensor features, $g$ is the active goal and $a^*$ is the teacher action. Before target binding, $g$ represents a startup or inspection skill. Deployed actors use measured inputs and the active goal; they cannot access true pose, teacher actions or a pending reply. The execution actor also receives the goal, whose separate benefit is not isolated in our experiments. Controller fitting holds the student fixed, and waiting-actor updates leave the execution actor and stopping rule unchanged.

\subsection{RGB perception and closed-loop deployment}
For cube control, a frozen CNN takes a $120\times120$ RGB crop with hand occlusion and predicts image coordinates and depth for the object center and eight keypoints. Camera calibration and rigid fitting recover position, XYZW orientation and a fitting error. Following ViserDex, Gaussian-splat rendering supplies training images for this estimator~\cite{bhardwaj2026viserdex}. The deployed model estimates pose from camera RGB.

At each control step, we estimate pose, run the recurrent student and local controller, and send joint targets. New agent observations reflect motion during inference and execution. All four physical-object systems use this RGB-to-pose-to-control interface. In simulation, cube uses RGB pose, while other objects use noisy simulator-derived pose. True pose supplies training supervision and independent scoring. Stopping and scoring are defined in Sec.~\ref{sec:setup}.
